% sec-03: the pharmacologic case and the simulation record.
\section{Pharmacologic Case}

Daraxonrasib (RMC-6236) is a RAS(ON) multi-selective inhibitor. RASolute 302
reported median overall survival of \textbf{13.2 months} against \textbf{6.6}
for chemotherapy in the RAS G12 population of previously treated metastatic
pancreatic cancer \cite{rasolute302}. A ten-arm quantitative systems
pharmacology simulation with 250 ordinary differential equations per patient had
returned \textbf{12.8} against \textbf{5.4 months} nine months earlier
\cite{qspmetpancre}. That proximity is chronology, not validation: 1000
simulated against 241 enrolled, a combination against a single agent, KRAS G12C
against a primarily G12D and G12V population.

\begin{apptable}
\begin{tabularx}{\textwidth}{>{\raggedright\arraybackslash}p{3.15cm} >{\raggedright\arraybackslash}p{1.85cm} >{\raggedright\arraybackslash}p{1.85cm} >{\raggedright\arraybackslash}X}
\headrow \thead{Source and date} & \thead{Test arm} & \thead{Control} & \thead{Reading and limitation}\\
\midrule
QSP simulation, Aug 2025 & 12.8 mo mOS & 5.4 mo mOS & Best arm HR 0.25; assumes no acquired resistance and ideal pharmacodynamics \\
Digital twin, Oct 2025 & 12.1 mo mOS & not applicable & Best mPFS HR 0.31; no patient-specific PK or tumour growth parameters \\
RASolute 302, May 2026 & 13.2 mo mOS & 6.6 mo mOS & Metastatic, previously treated; says nothing about the resectable setting \\
\end{tabularx}
\end{apptable}
\tabcap{Each source with its stated limitation in the same row, so no line of the
table can be read without the caveat that qualifies it. Limitations are the
authors' own, not added here.}

Neither simulation nor the metastatic readout establishes anything about
perioperative exposure in resectable disease, which is what this concept
proposes to measure directly in the resection specimen.
